\section{Preliminaries}

A finite normal-form game is $\Gamma=(N,(A_i)_{i\in N},(u_i)_{i\in N})$, where $N=\{1,\ldots,n\}$, $A_i$ is player $i$'s finite action set, $A=\prod_i A_i$, and $u_i:A\to\mathbb R$.  We write $a=(a_i,a_{-i})$ for a pure profile and $\sigma=(\sigma_i)_{i\in N}$ for an independent mixed profile.  Player $i$'s regret is
\[
\operatorname{Reg}_i^\Gamma(\sigma)
=
\max_{\tau_i\in\Delta(A_i)}
\bigl[u_i(\tau_i,\sigma_{-i})-u_i(\sigma)\bigr].
\]
The profile is an $\epsilon$-Nash equilibrium if every regret is at most $\epsilon$.

\paragraph{Relabelings and exact symmetries.}
A relabeling $g$ consists of a permutation $\pi_g$ of the players and, for every player $i$, a bijection
$\tau_{g,i}:A_i\to A_{\pi_g(i)}$.  It maps a pure profile according to
\[
(ga)_{\pi_g(i)}=\tau_{g,i}(a_i).
\]
The induced action on games is defined by
\[
(g\Gamma)_{\pi_g(i)}(ga)=u_i(a).
\]
Thus $g$ is an exact symmetry of $\Gamma$ precisely when $g\Gamma=\Gamma$.  We consider a finite group $G$ of such relabelings, represented algorithmically by a generating set $S$.  A mixed profile respects $G$ if
\[
\sigma_{\pi_g(i)}\bigl(\tau_{g,i}(a_i)\bigr)=\sigma_i(a_i)
\quad
\text{for all }g\in G,i,a_i.
\]
Equivalently, strategy probabilities are constant on the orbits of the disjoint union $\bigsqcup_i A_i$.

\paragraph{Strategic distance.}
For compatible games $\Gamma$ and $\Gamma'$ with utilities $u$ and $v$, define
\begin{align}
\devdist(\Gamma,\Gamma')
=\max_{\substack{i,a_{-i},\\a_i,b_i\in A_i}}
\Big|&[u_i(a_i,a_{-i})-u_i(b_i,a_{-i})] \nonumber\\
&-[v_i(a_i,a_{-i})-v_i(b_i,a_{-i})]\Big|.
\label{eq:mpd}
\end{align}
This is the maximum pairwise difference of unilateral incentives.  If
$r_i(a)=u_i(a)-v_i(a)$ is the residual, then
\begin{equation}
\devdist(\Gamma,\Gamma')
=
\max_{i,a_{-i}}
\left(
\max_{a_i}r_i(a_i,a_{-i})-\min_{a_i}r_i(a_i,a_{-i})
\right).
\label{eq:range-form}
\end{equation}

A game component $q$ is \emph{nonstrategic} when
$q_i(a_i,a_{-i})=c_i(a_{-i})$ for some functions $c_i$; such a component changes payoff levels but no unilateral incentive.

\begin{proposition}[Strategic quotient norm]
\label{prop:quotient}
The function $\devdist$ is a relabeling-invariant seminorm distance, and
\[
\devdist(\Gamma,\Gamma')=0
\quad\Longleftrightarrow\quad
\Gamma-\Gamma'\text{ is nonstrategic}.
\]
Consequently it is a norm on the quotient of games by the nonstrategic subspace.
\end{proposition}
\begin{proof}
It is the maximum of absolute linear functionals, which gives symmetry, homogeneity, and the triangle inequality.  Distance zero means that, for every fixed $(i,a_{-i})$, the residual $r_i(a_i,a_{-i})$ is constant over $a_i$, exactly the nonstrategic form.  Relabelings only permute the functionals in the maximum.
\end{proof}

For comparison, let
$\|\Gamma-\Gamma'\|_\infty=\max_{i,a}|u_i(a)-v_i(a)|$.
The triangle inequality for real numbers gives
\begin{equation}
\devdist(\Gamma,\Gamma')
\le 2\|\Gamma-\Gamma'\|_\infty.
\label{eq:raw-upper}
\end{equation}
There is no reverse inequality up to any universal constant: adding an arbitrarily large nonstrategic component leaves $\devdist$ equal to zero.  This separation is why we define almost symmetry through incentives rather than raw utilities.
